% 04_macros.tex
% Componentes Especiales, Avisos y Formatos Personalizados

% 1. Formato para Comandos y Rutas (Courier, 10pt)
\newcommand{\comando}[1]{{\fontfamily{pcr}\fontsize{10}{12}\selectfont\texttt{#1}}}

% 2. Formato para Resultados Esperados (Verde)
\newcommand{\resultado}[1]{\textcolor{colorSecundario}{\textbf{#1}}}

% 3. Componentes Especiales de Usabilidad (Estándar ISO)

% Nota: Consejos de productividad (Fondo azul claro, icono info)
\newtcolorbox{AvisoNota}{
    colback=bgNota,
    colframe=bgNota!80!black, % Borde ligeramente más oscuro
    boxrule=0.5pt,
    leftrule=4pt,
    arc=2pt,
    title={\faInfoCircle\ \textbf{Nota}},
    coltitle=black,
    fonttitle=\sffamily,
    attach title to upper=\quad
}

% Precaución: Evitar errores operacionales (Borde amarillo, texto negrita, sin fondo invasivo)
\newtcolorbox{AvisoPrecaucion}{
    colback=white,
    colframe=colorAcento,
    boxrule=1.5pt,
    arc=2pt,
    fontupper=\bfseries,
    title={\faExclamationTriangle\ \textbf{Precaución}},
    coltitle=black,
    fonttitle=\sffamily,
    attach title to upper=\par
}

% Advertencia: Prevenir pérdida de datos (Fondo rojo suave, icono peligro)
\newtcolorbox{AvisoAdvertencia}{
    colback=bgAdvertencia,
    colframe=red!70!black,
    boxrule=0.5pt,
    leftrule=4pt,
    arc=2pt,
    title={\faRadiation\ \textbf{Advertencia}},
    coltitle=black,
    fonttitle=\sffamily,
    attach title to upper=\par
}

% 4. Formato de Checklists de Verificación
% Utiliza cuadros de marcado estandarizados
\newlist{checklist}{itemize}{1}
\setlist[checklist]{label=$\square$, leftmargin=*}